\section{引言}
\label{sec:intro}

慢故障是分散式系統中最隱蔽的故障類型之一~\cite{huang2017gray, gunawi2018fail}。
在慢故障中，節點仍能回應請求，但處理效能顯著下降。
崩潰故障可透過健康檢查立即偵測，慢故障則不然：
延遲升高可能源自暫態負載而非真正的效能降級，
故障訊號因此模糊不清。
若未及時處置，32 節點叢集中的單一慢節點可使系統 P99 延遲增加
3--9$\times$，違反服務等級目標（SLO），
並沿服務依賴鏈蔓延至整個系統~\cite{dean2013tail}。

處理尾延遲的標準做法是\emph{對沖請求}（hedged requests）：
等待一小段時間後，將同一請求的副本發送至另一台伺服器，
取最先到達的回應~\cite{dean2013tail}。
對沖請求能有效處理\emph{隨機慢請求}——
即垃圾回收、上下文切換或背景 I/O 造成的暫態延遲。
但持續性慢故障與此截然不同：
同一節點在數千個請求中持續緩慢，
逐次請求的對沖機制無法解決問題，
因為後續請求仍會持續路由至同一故障節點。

現代負載均衡器主要使用 Power-of-Two-Choices（P2C）
路由~\cite{mitzenmacher2001power}，
依權重隨機抽取兩個候選工作節點，
並將請求路由至其中佇列較短者。
由此衍生一個核心問題：\emph{P2C 的佇列深度感知
是否已能處理慢故障，使顯式緩解成為多餘？}
本文以三個理論結果及廣泛的實證加以解答。

\begin{enumerate}
\item \textbf{操作包絡線}：任何緩解策略可容忍的最大故障節點比例為
  \begin{equation}
    \fmax = 1 - \frac{L}{U}
    \label{eq:envelope}
  \end{equation}
  其中 $L$ 為系統負載因子，
  $U \approx 0.89$ 為 M/D/1 排隊延遲開始超線性增長的利用率上限。
  此公式毋須擬合參數，即能預測四個負載等級的成功／失敗邊界
  （第~\ref{sec:envelope}~節）。

\item \textbf{嚴重度非單調性}：在 P2C 路由下，
  「最糟」的減速因子並非最高值，
  而是某個特定值。在此減速因子下，
  到達故障節點的流量比例恰好等於目標百分位的尾端寬度：
  \begin{equation}
    \sstar = \frac{f \cdot q}{(1-q)(N-f)}
    \label{eq:sstar}
  \end{equation}
  高於 $\sstar$ 時，P2C 透過佇列深度比較自動將流量導離故障節點，
  系統無需顯式介入即可\emph{自行隔離}故障節點
  （第~\ref{sec:nonmono}~節）。

\item \textbf{SHED--P2C 互補性}：P2C 路由與顯式權重削減（SHED）
  作用於路由決策的不同階段：
  SHED 控制\emph{候選者被選取的機率}，
  而 P2C 提供\emph{佇列深度的決勝機制}。
  兩者互補而非冗餘。
  SHED 的效益隨 $O(N^{-\alpha})$（$\alpha \approx 1.3$）衰減，
  在中度嚴重度（$2$--$5\times$）時最大化，
  而在 $\sstar$ 以上（P2C 自行隔離之處）則趨於消失。
  關鍵的是，在 P2C 下最低權重 $w_{\min} = 0.05$
  可達到與完全隔離（$w = 0$）幾乎相同的 P99，
  證明\emph{偵測精確度遠不如偵測速度重要}
  （第~\ref{sec:complementarity}~節）。
\end{enumerate}

上述結果經 8 至 64 節點叢集的離散事件模擬加以驗證，
涵蓋 15 個故障場景。評估比較八種緩解策略，
包括三種文獻基線（對沖請求~\cite{dean2013tail}、
異常值黑名單~\cite{lakshman2010cassandra}和超時重試）。
主要實證發現如下：

\begin{itemize}
\item 偵測速度是首要瓶頸：在 P2C 路由下，
  漸進式權重削減與完全隔離產生幾乎相同的 P99 延遲
  （差距 $< 2$\,ms）。一旦偵測到故障節點，
  回應策略便無關緊要。
  此結論在不同叢集大小、嚴重度等級與負載因子下皆成立
  （第~\ref{sec:eval-shed}~節）。

\item 反應式機制失效：對沖請求和超時重試
  與無緩解的情況相比，SLO 達成率僅提升不到 0.3 個百分點——
  它們缺乏偵測機制，是為隨機慢請求而非持續故障設計的
  （第~\ref{sec:eval-lit}~節）。

\item 主動偵測既必要又充分：所有能偵測故障節點
  並重新導向流量的策略，在 80\% 負載下，
  於所有場景中均達成 P99 $< 45$\,ms。
  無緩解與任何主動策略的差距（5--70\,ms）
  遠超不同主動策略之間的差距
  （第~\ref{sec:eval-tiers}~節）。
\end{itemize}

本文的實務貢獻如下：
（1）判斷給定負載與故障比例下緩解是否可行的簡明公式；
（2）識別最危險減速區間的分析工具；
（3）透過互補性定理釐清何時 P2C 本身即已足夠、何時需要顯式偵測；
（4）證明偵測速度（而非回應精確度）才是慢故障緩解系統首要設計目標的實證依據。
